\chapter*{Contribuciones Personales}
\label{cap:contribucionesPersonales}
\addcontentsline{toc}{chapter}{Contribuciones Personales}

En caso de trabajos no unipersonales, cada participante indicará en la memoria su contribución al proyecto con una extensión de al menos dos páginas por cada uno de los participantes.

En caso de trabajo unipersonal, elimina esta página en el fichero \texttt{TFGTeXiS.tex} (comenta o borra la línea \verb|\include{Capitulos/ContribucionesPersonales}|).

\section*{Iván Alcalde Cámara}
Al menos dos páginas con las contribuciones.

\section*{Contribuciones de Ignacio Melendo Bruggeman}

En el desarrollo de la aplicación \textit{ClimbIt} ha sido posible gracias al esfuerzo colaborativo junto a Iván Alcalde. Con el objetivo de aumentar el alcance de la aplicación lo máximo posible sin comprometer la calidad, se tomo la decisión de hacer una división de responsabilidades. A continuación se destallan mis contribuciones en las diferentes areas del proeycto.

\subsection*{1. Gestión del Proyecto y Metodología}

En el inicio del proyecto me encargué de la configuración de el marco metodológico, preparando el entorno de \textit{GitHub Project} para su posterior uso con la metodología definida de \textit{Scrumban}. También establecí politicas del de control de versiónes en el repositorio de GitHub, configurando sistemas de protección de ramas en \textit{main y develop} y mecanismos para la integración de código.

Durante la fase de definición del valor del producto, elaboré el \textit{Value Proposition Canvas} para empatizar con los usuarios finales definiendo de esta manera los \textit{Pains}, \textit{Gains} que encuentran los escaladores en la actualidad. Una vez definido y utilizando los resultados obtenidos generé un \textit{User Story Map} inicial que fue imprescindible para visualizar el viaje del usuario y definir los hitos del proyecto como el Producto Mínimo Viable (MVP). También me encargué de la redacción y mantenimiento de documentos técnicos del repositorio como el Readme para la guía de instalación y el Changelog para mantener un historial del versionado siguiendo los estadares de \textit{keep a changelog}.

\subsection*{2. Análisis de Requisitos y Diseño del Sistema}

En la fase de análisis de requisitos mi aportación se centro en aplicar el Diseño Centrado en el Usuario (DCU) mediante sistemas como la definición de \textit{Personas} como la escaladora Klaudia y el \textit{route-setter} Victor (gestor del rocódromo). Basandonos en estos perfiles y el USM junto con mi compañero fuimos definiendo las Historias de Usuario de manera progresiva a lo largo de los sprints. Una vez alcanzados los hitos definidos en el Mapa de Historia de Usuario me encargue de la publicación de las releases en la plataforma GitHub.

En cuanto al diseño del sistema, participé activamente en la toma de decisiones arquitectónicas. Contribuí al diseño de la arquitectura hexagonal para el \textit{backend} y la arquitectura modular basada en capas para el \textit{frontend}. También colaboré en el diseño del modelo entidad-relación de la base de datos PostgreSQL, asegurando la correcta normalización y estructuración de la información. 

En cuando al diseño de la arquitectura del sistema, junto con mi compañero fuimos tomando las decisiones técnicas para el uso de la arquitura hexagonal para el \textit{backend} y la arquitectura por módulos del \textit{frontend}. Además colaboré en el diseño inicial del modelo de entidad-relación de la base de datos PostgreSQL definiendo todas las entidades y atributos principales.

En el tópico del diseño de interfaces y experiencia de usuario (\textit{UI/UX}) inicialmente generé una estructura de navegación de las diferentes interfaces del sistema mediante bocetos como son las interfaces de registro, y de las entidades del perfil, rocódromos zonas y rutas. Posteriormente durante la fase de desarrollo del proyecto junto con mi compañero fuimos adaptando y creando las vistas necesarias. Entre estas vistas destaco la interfaz de información de la zona la cual fue una adaptación de los sistemas de \textit{TopLogger y WeClimb} la cual requirió de un sistema para generar un mapa interactivo y un sistema de tarjetas para las diferentes rutas de un rocódromo.

\subsection*{3. Implementación y Desarrollo Full-Stack}
Durante la fase de codificación, ambos autores adoptamos un enfoque de trabajo \textit{Full-Stack}, interviniendo en todas las capas del sistema según las necesidades del \textit{sprint}. Sin embargo, mi carga de trabajo técnico estuvo focalizada de manera principal en el desarrollo del \textit{frontend}. Me encargué de materializar los diseños UI/UX previamente definidos, implementando una \textit{Single Page Application} (SPA) robusta, desarrollando el cliente HTTP centralizado para el consumo de la API y garantizando que la interfaz fuera completamente responsiva y adaptada a dispositivos móviles.

Durante el desarrollo propio de \textit{ClimbIt} tanto mi compañero como yo intervinimos en mayor o menor medida en todas las capas del proyecto, sin embargo creamos una separación de responsabilidades en la cual centré más mi atención en el desarrollo del frontend.

\subsection*{4. Integración Continua, Despliegue y DevOps (CI/CD)}
Como responsable de la infraestructura y automatización, diseñé y ejecuté toda la arquitectura de despliegue del proyecto. Creé los \textit{pipelines} de Integración Continua (CI) mediante \textit{GitHub Actions} para los entornos de preproducción y producción. Fui el encargado de configurar físicamente el servidor \textit{self-hosted} utilizando una Raspberry Pi 4, donde sistematicé el despliegue de los contenedores \textit{Docker}.

Para garantizar la seguridad y accesibilidad del sistema desde el exterior, gestioné la configuración de los dominios y establecí el enrutado seguro mediante túneles \textit{Cloudflare Zero Trust}. Finalmente, asumí el reto técnico de empaquetar la aplicación web progresiva (PWA) utilizando la tecnología \textit{Trusted Web Activity} (TWA) y \textit{Bubblewrap}, logrando compilar la versión nativa y gestionando exitosamente su subida, verificación y publicación en la \textit{Google Play Console}.

\subsection*{5. Pruebas, Validación y Trabajo de Campo}
Para asegurar que la aplicación aportaba un valor real al sector, lideré el contacto directo y la comunicación técnica con diversos rocódromos y escaladores de la comunidad. Realicé un extenso trabajo de campo que incluyó la creación de mapas vectoriales interactivos para las instalaciones y el mapeo manual de las rutas en los rocódromos colaboradores, destacando el trabajo realizado con el centro UADIBLOC y las instalaciones de Hoyo de Manzanares.

Asimismo, coordiné y ejecuté las sesiones de pruebas con usuarios finales (User Testing). Estas sesiones se dividieron en dos grandes fases: una evaluación intermedia durante el desarrollo para ajustar la usabilidad de las vistas, y unas pruebas exhaustivas al final del desarrollo para validar el MVP y el sistema de gamificación.

\subsection*{6. Redacción de la Memoria}
Por último, en lo que respecta a la elaboración de este documento académico, he sido el autor principal y responsable de la redacción, estructuración y formato LaTeX de los siguientes capítulos:
\begin{itemize}
    \item Resumen y Abstract.
    \item Capítulo 1: Introducción.
    \item Capítulo 2: Estado del Arte.
    \item Capítulo 3: Metodología de Desarrollo y Planificación.
    \item Capítulo 4: Análisis y Especificación de Requisitos.
    \item Capítulo 7: Integración Continua, Despliegue y DevOps.
    \item Capítulo 9: Conclusiones y Trabajo Futuro.
\end{itemize}